\scanpage{49}\mspage{36}

pour la conclusion du lemme, inutile de supposer $C$ stable par compositions infinies.

Prouvons 2.1. $3^\circ$. On a $a\Longrightarrow b$, car par 2.1. $1^\circ$\unsure{,|)} $u_\bullet\in(TC)_I$ implique que pour $i\in I$, $u_i:X_i\longrightarrow Y_i$ est dans $TC$, \deleted{donc dans} (cf cor.\ 2.1. \unsure{a}), \unsure{Donc} il reste à prouver $\mathrm{b)}\Longrightarrow\mathrm{a)}$ i.e. si $u_\bullet\in C\cap W_{I^{\unsure{o|\circ}}}$, alors $u_\bullet\in(TC)_I$, i.e. les $u_{K,I_{/i_0}}$ sont dans $TC$ \deleted{\unsure{puisq.}}

On sait déjà que $u_{K,I_{/i_0}}\in C$ par hyp. $u_\bullet\in C_I$, donc il faut prouver que c'est $\in W$. On est ramené à prouver que $u_K\in W$, i.e. $u_K\in TC$, car \deleted{par} le diagramme (où $J=I_{/i_0}$)

\diagram{\begin{tikzcd}[ampersand replacement=\&,column sep=huge]
 X_K\arrow[r]\arrow[d,"u_K"']\arrow[dr,phantom,"\text{cocart}"] \& X_J\arrow[d,"\widetilde u_{K,J}"]\arrow[dd,bend left=60,"u_J"]\\
 Y_{K.}\arrow[r,dash] \& Y_{K,J}\arrow[d,"u_{K,J}"]\\
 \&Y_J
\end{tikzcd}}

$u_K\in TC$ implique $\widetilde u_{K,J}\in TC$, et comme $u_J\simeq u_{i_0}\in W$ par hypothèse, il en résulte bien $u_{K\unsure{J|,J}}\in W$ par saturation. \inserted{Mais on va}

\deleted{Donc on est ramené à prouver le \uline{Lemme} 2.9. \inserted{Supposons $I$ artinien.} Soit $u_\bullet\in C_I\cap W_{I^o}$. Alors pour \unsure{tt} ouvert $K$ de $I$, $u_K:X_K\longrightarrow Y_K$ est dans $\deleted{\unsure{T}}TC$ (donc dans $W$),} s'y prendre autrement:

% archive sheet number: 48

